%\documentclass[12pt,oneside]{amsart}
\documentclass{scrartcl}

%\documentclass{article}

\usepackage{amsthm,amsmath,amssymb}
\usepackage[numbers]{natbib}

%\usepackage[colorlinks]{hyperref}
\usepackage{hyperref}
\hypersetup{
    colorlinks,%
    citecolor=black,%
    filecolor=black,%
    linkcolor=black,%
    urlcolor=black
}
\usepackage{hypernat}
\usepackage[top=2.5cm, bottom=2.5cm, left=2.8cm,
right=2.8cm]{geometry}
\usepackage{graphicx}


%\setlength{\parskip}{0pt}
%\setlength{\parsep}{0pt}
%\setlength{\headsep}{0pt}
%\setlength{\topskip}{0pt}
%\setlength{\topmargin}{0pt}
%\setlength{\topsep}{0pt}
%\setlength{\partopsep}{0pt}

\linespread{1}

%\usepackage{marvosym}

%\textwidth = 6.5 in \textheight = 9.2 in \oddsidemargin = 0.0 in
%\evensidemargin = 0.0 in \topmargin = -0.0 in \headheight = 0.0 in
%\headsep = 0.0in \parskip = 0.0in \parindent = 0.0in
%\def\baselinestretch{0.871}


\newtheorem{theorem}{Theorem}[section]
\newtheorem{proposition}[theorem]{Proposition}
\newtheorem{problem}[theorem]{Problem}
\newtheorem{fact}[theorem]{Fact}
\newtheorem{lemma}[theorem]{Lemma}
\newtheorem{defn}[theorem]{Definition}
\newtheorem{corollary}[theorem]{Corollary}
\newtheorem{remark}{Remark}[section]
\newtheorem{example}[theorem]{Example}

\newcommand{\E}{{\mathbb E}}
\newcommand{\se}{{\mathcal{E}}}
\newcommand{\R}{{\mathbb R}}
\renewcommand{\P}{{\mathbb P}}
\newcommand{\ac}{{\mathcal{A}}}
\newcommand{\A}{{\mathcal{A}}}
\newcommand{\B}{{\mathcal{B}}}
\newcommand{\C}{{\mathcal{C}}}
\newcommand{\M}{{\mathcal{M}}}
\newcommand{\Mf}{{\mathfrak{M}}}
\newcommand{\T}{{\mathcal{T}}}
\newcommand{\Z}{{\mathcal{Z}}}
\newcommand{\K}{{\mathcal{K}}}
\newcommand{\lc}{{\mathcal{L}}}
\newcommand{\Ps}{{\mathcal{P}}}
\newcommand{\G}{{\mathcal{G}}}
\newcommand{\pa}{{\mathring{p}}}
\newcommand{\F}{{\cal F}}
\newcommand{\samp}{{\mathcal{X}}}
\newcommand{\X}{{\mathcal{X}}}
\newcommand{\hi}{{\hat{\Phi}}}
\newcommand{\data}{{\text{data}}}

\newcounter{rcnt}[section]
\renewcommand{\thercnt}{(\roman{rcnt})}

\def\qt#1{\qquad\text{#1}}

\def\argmin{\mathop{\rm argmin}}
\def\argmax{\mathop{\rm argmax}}


\setlength{\parskip}{1.4 \medskipamount}
%\sloppy
%\linespread{1.3}

\begin{document}

\title{STAT 238 - Bayesian Statistics  \\
  Lecture Thirteen} 
\subtitle{Spring 2026, UC Berkeley}

\author{Aditya Guntuboyina}


\date{20 Feb 2026}

\maketitle


\section{Dirichlet Prior and Multinomial Likelihood}

We saw the following fact in the last lecture: 
\begin{align*}
&  (p_1, \dots, p_k) \sim \text{Dirichlet}(a_1, \dots, a_k)  \text{
  and }  (X_1, \dots, X_k) \mid (p_1, \dots, p_k) \sim 
  \text{Multinomial}(n; p_1, \dots, p_k) \\
& \implies (p_1, \dots, p_k) \mid X_1 = x_1, \dots, X_k = x_k \sim
  \text{Dirichlet}(a_1 + x_1, a_2 + x_2, \dots, a_k + x_k). 
\end{align*}
Here the data are given by $x_1, \dots, x_k$ where $x_i$ represents
count of the $i$-th outcome in $n$ repeated trials (so $x_1 + \dots + x_k =
n$). The probabilities $p_1, \dots, p_k$ associated with $k$
outcomes represent the unknown parameters which are to
be estimated. 

The posterior mean estimate of $p_i$ is given by:
\begin{align*}
  \E (p_i \mid X_1 = x_1, \dots, X_k = x_k) = \frac{a_i + x_i}{(a_1 +
  x_1) + \dots + (a_k + x_k)} = \frac{a_i + x_i}{n + a_1 + \dots +
  a_k}
\end{align*}
where we used $x_1 + \dots + x_k = n$.

\section{Special Case: the $\text{Dirichlet}(0, \dots, 0)$ prior}

In the special case where all $a_i = 0$, then the posterior mean of
$p_i$ is simply $x_i/n$; so we are estimating the probability of the
$i$-th outcome by simply the proportion of times the $i$-th outcome
appeared in the $n$ trials. This is just the frequentist MLE. So the
posterior mean estimate coincides with the frequentist MLE when the
prior is $\text{Dirichlet}(0, \dots, 0)$.

More precisely, the posterior for the
$\text{Dirichlet}(0, \dots, 0)$ prior is $\text{Dirichlet}(x_1,
\dots, x_k)$. A key property of this posterior is that if $x_i = 0$ for
some $i$, then the posterior
places all its mass on the set $\{p_i = 0\}$. In other words, any
outcome that is not observed in the sample is assigned zero
probability under the posterior and is therefore completely excluded
from future consideration.

As an explicit example, suppose $k = 6$, $n = 4$, $x_1 = 0, x_2 = 0,
x_3 = 0, x_4 = 0, x_5 = 3, x_6 = 1$. In other words, we rolled a
six-faced die $4$ times; in 3 of the 4 rolls, we observed the outcome
5, and in the other roll, we observed a 6. After observing this data
(note that we started with the prior $\text{Dirichlet}(0, 0, 0, 0, 0,
0)$), our posterior becomes $\text{Dirichlet}(0, 0, 0, 0, 3, 1)$. This
means that the posterior is concentrated on the probabilities $(0, 0,
0, 0, p_5, 1 - p_5)$ with $p_5 \sim \text{Beta}(3, 1)$.

\section{Further special case: $\text{Dirichlet}(0, \dots, 0)$ prior
  and data in which no outcome is repeated}

Suppose the actual outcomes observed in the $n$ trials are $y_1,
\dots, y_n$ and assume that these are all distinct (which means that
we did not observe any repeated outcome). As before, we work with the
$\text{Dirichlet}(0, \dots, 0)$ prior. Then the posterior can be
written as:
\begin{align*}
  \sum_{i=1}^n w_i \delta_{\{y_i\}} ~~ \text{ where } (w_1, \dots,
  w_n) \sim \text{Dirichlet}(1, \dots, 1), 
\end{align*}
where $\delta_{\{y_i\}}$ denotes the point mass at the point $y_i$. 

Note also that the density corresponding to the $\text{Dirichlet}(1,
\dots, 1)$ distribution is constant, so the above posterior can be
understood as the uniform distribution over all probability measures
that are supported on the observed data points.

Posterior samples therefore can be generated by:
\begin{align*}
  (w_1^{(j)}, \dots, w_n^{(j)}) \overset{\text{i.i.d}}{\sim}
  \text{Dirichlet}(1, \dots, 1)  
\end{align*}
for $j = 1, \dots, M$ (for a large number $M$). Inference done with
these samples is referred to as the \textbf{Bayesian Bootstrap} (see
\citet{Rubin1981BayesianBootstrap}). This is because the generation
process for these samples is similar to the usual nonparametric bootstrap
sample generation. Indeed, the usual bootstrap samples 
can be seen as being generated from:
\begin{align*}
  \sum_{i=1}^n v_i \delta_{\{y_i\}} ~~ \text{ where } (v_1, \dots,
  v_n) \sim \frac{1}{n} \text{Multinomial}(n; 1/n, \dots, 1/n). 
\end{align*}
The distributions $(w_1, \dots, w_n) \sim \text{Dir}(1,
\dots, 1)$ $\&$ $(v_1, \dots, v_n) \sim \text{Mult}(n; 1/n, \dots,
1/n)/n$ are quite close when $n$ is large. For example,
\begin{align*}
  & \E w_i = \E v_i = 1/n \\
  & \text{var}(w_i) = \frac{n-1}{n^2(n+1)} \approx \text{var}(v_i) =
    \frac{n-1}{n^3} \\
  & \text{correlation}(w_i, w_j) = \text{correlation}(v_i, v_j) =
    \frac{-1}{n-1} ~~ \text{ when } i \neq j. 
\end{align*}
So operationally, bootstrap and the Bayesian bootstrap give very
similar results, although conceptually they are quite different. For
more information on the Bayesian bootstrap and its relation to the
classical bootstrap, see the blog post by Rasmus B{\aa}{\aa}th:
\url{https://www.sumsar.net/blog/2015/04/the-non-parametric-bootstrap-as-a-bayesian-model/}.    



\begin{thebibliography}{99}

\bibitem[Rubin(1981)]{Rubin1981BayesianBootstrap}
D.~B. Rubin,
\newblock ``The Bayesian bootstrap,''
\newblock \emph{The Annals of Statistics}, vol.~9,
no.~1, pp.~130--134, 1981.
  
%\bibitem[Fisher(1934)]{Fisher1934}
%R.~A. Fisher,
%\newblock ``Probability, likelihood and quantity of information in the logic of
%uncertain inference,''
%\newblock \emph{Proceedings of the Royal Society of London. Series~A,
%Containing Papers of a Mathematical and Physical Character}, vol.~146,
%no.~856, pp.~1--8, 1934.

%\bibitem[Efron(2010)]{Efron}
%Efron, B. (2010)
%\textit{Large-scale inference: empirical Bayes methods for estimation,
%testing, and prediction}. Cambridge University Press, 2010.
  
%     \bibitem[Jaynes(2003)]{Jaynes} Jaynes, E. T, (2003)
%  \textit{Probability theory: the logic of science}. Cambridge
%  University Press, 2003.

%  \bibitem[Sinai(1992)]{Sinai} Sinai, Y. G, (1992)
%  \textit{Probability theory: an introductory course}. Springer
%  Verlag, 1992.  

%\bibitem[{deGroot(1986)}]{DeGrootBlackwell}DeGroot, M. H. (1986)
%       A conversation with David Blackwell.
%\newblock \emph{Statistical Science}, \textbf{1}, 40--53.

%         \bibitem[Mosteller(1987)]{Mosteller} Mosteller, F, (1987)
%  \textit{Fifty challenging problems in probability with
%    solutions}. Courier Corporation, 1987.

%    \bibitem[MacKay(2003)]{MacKay} MacKay, D. J. C., (2003)
%  \textit{Information theory, inference, and learning
%  algorithms}. Cambridge University Press, 2003.

%\bibitem[{Lindley(2013)}]{Lindley}Lindley, D. V. (2013)
%   \textit{Understanding Uncertainty}. John Wiley and sons, 2013.


\end{thebibliography}







\end{document}

%%% Local Variables:
%%% mode: latex
%%% TeX-master: t
%%% End:
